\documentclass{article}

\usepackage{amsmath}
\usepackage{amssymb}
\usepackage{geometry}
\geometry{margin=1in}

\title{Deep ML Techniques}
\author{Mobina Amrollahi}
\date{\today}

\begin{document}

\maketitle

\tableofcontents

This document covers machine learning techniques explored in the deep-ml problem set.

\section{Probability and Statistics}
\label{sec:prob-stat}

\subsection{Distributions}
\label{subsec:prob-stat:dist}

\subsubsection{Exponential}
\label{subsubsec:prob-stat:dist:exp}

\paragraph{PDF}
\label{para:prob-stat:dist:exp:pdf}
\begin{equation}
    \label{eq:exp-pdf}
    p(x; \lambda) =
    \begin{cases}
        \lambda e^{-\lambda x}, & x \geq 0 \\
        0, & x < 0
    \end{cases}
\end{equation}

\paragraph{CDF}
\label{para:prob-stat:dist:exp:cdf}
\begin{equation}
    \label{eq:exp-cdf}
    P(x; \lambda) = P(X \leq x) =
    \begin{cases}
        1 - e^{-\lambda x}, & x \geq 0 \\
        0, & x < 0
    \end{cases}
\end{equation}

\paragraph{Mean}
\label{para:prob-stat:dist:exp:mean}
\begin{equation}
    \label{eq:exp-mean}
    \mathbb{E}[X] = \frac{1}{\lambda}
\end{equation}

\paragraph{Variance}
\label{para:prob-stat:dist:exp:variance}
\begin{equation}
    \label{eq:exp-variance}
    \operatorname{Var}(X) = \frac{1}{\lambda^2}
\end{equation}

\subsubsection{Normal (Gaussian)}
\label{subsubsec:prob-stat:dist:norm}

\paragraph{PDF}
\label{para:prob-stat:dist:norm:pdf}
\begin{equation}
    \label{eq:norm-pdf}
    p(x;\, \mu, \sigma^2) = \frac{1}{\sigma\sqrt{2\pi}}\exp\!\left(-\frac{(x-\mu)^2}{2\sigma^2}\right)
\end{equation}

\paragraph{CDF}
\label{para:prob-stat:dist:norm:cdf}
\begin{equation}
    \label{eq:norm-cdf}
    \Phi(x;\, \mu, \sigma^2) = P(X \leq x) = \frac{1}{2}\left[1 + \operatorname{erf}\!\left(\frac{x - \mu}{\sigma\sqrt{2}}\right)\right]
\end{equation}

\paragraph{Mean}
\label{para:prob-stat:dist:norm:mean}
\begin{equation}
    \label{eq:norm-mean}
    \mathbb{E}[X] = \mu
\end{equation}

\paragraph{Variance}
\label{para:prob-stat:dist:norm:variance}
\begin{equation}
    \label{eq:norm-variance}
    \operatorname{Var}(X) = \sigma^2
\end{equation}

\subsubsection{Bernoulli}
\label{subsubsec:prob-stat:dist:bern}

\paragraph{PMF}
\label{para:prob-stat:dist:bern:pmf}
\begin{equation}
    \label{eq:bern-pmf}
    P(X = x) =
    \begin{cases}
        1 - p, & x = 0 \\
        p, & x = 1
    \end{cases}
\end{equation}

\paragraph{CDF}
\label{para:prob-stat:dist:bern:cdf}
\begin{equation}
    \label{eq:bern-cdf}
    P(X \le x) =
    \begin{cases}
        0, & x < 0 \\
        1-p, & 0 \le x < 1 \\
        1, & x \ge 1
    \end{cases}
\end{equation}

\paragraph{Mean}\label{para:prob-stat:dist:bern:mean}
\begin{equation}
    \label{eq:bern-mean}
    \mathbb{E}[X] = 0 \times (1 - p) + 1 \times p = p
\end{equation}

\paragraph{Variance}
\label{para:prob-stat:dist:bern:variance}
\begin{equation}
    \label{eq:bern-variance}
    \operatorname{Var}(X) = p (1-p)
\end{equation}

\subsection{Softmax}\label{subsec:prob-stat:softmax}
\subsubsection{Softmax Probabilities}
\begin{equation}
    \label{eq:softmax}
    \sigma(\mathbf{z}_i) = \frac{e^{z_i}}{\sum_{j=1}^{K} e^{z_j}}
\end{equation}

\subsubsection{Logarithm of Softmax Probabilities}\label{subsubsec:prob-stat:softmax:log}
\begin{equation}
    \label{eq:log-softmax}
    \log \sigma(\mathbf{z}_i) = z_i - \log \sum_{j=1}^{K} e^{z_j}
\end{equation}

\subsubsection{Stable Softmax}\label{subsubsec:prob-stat:softmax:stable}
\begin{equation}
    \label{eq:stable-softmax}
    \sigma(\mathbf{z}_i) = \frac{e^{z_i - \max{z_j}}}{\sum_{j=1}^{K} e^{z_j- \max{z_j}}}
\end{equation}

\subsection{Maximum Likelihood Estimation}
\label{subsec:prob-stat:maximum-likelihood-estimation}

\section{Linear Algebra}
\label{sec:linear-algebra}

\subsection{Transformation Matrices}
\label{subsec:linear-algebra:transformation-matrices}

\section{Numpy}
\label{sec:numpy}

\begin{itemize}
    \item \texttt{np.array(lst)} turns a list into a numpy array; \texttt{arr.tolist()} turns a numpy array into a list.
    \item \texttt{np.random.choice(elements, p=probabilities)} pick from a numpy array based on specific probabilities.
    \item \texttt{np.random.seed(seed)} sets the random seed for the purpose of reproducibility.
\end{itemize}

\end{document}
